% PAGES: 215-216
% Opens by closing the "\begin{proof}" of Lemma 7.3 left open at the end of
% sec07_c1_3.tex (folio 198), whose last item was equation (7.27). Do not open a
% new "\begin{proof}" for this continuation.

since the covariance matrix $\Sigma_{\bfX}$ is symmetric, i.e., $\Sigma_{\bfX}^t =
\Sigma_{\bfX}$.

Moreover, since $\left(C^{(1)}\right)^t\Sigma_{\bfX}C^{(2)} \in \R$ is a number, it
follows that
\begin{equation}
\left(\left(C^{(1)}\right)^t\Sigma_{\bfX}C^{(2)}\right)^t
\;=\; \left(C^{(1)}\right)^t\Sigma_{\bfX}C^{(2)}.
\end{equation}
From (7.27) and (7.28), we conclude that
\[
\left(C^{(1)}\right)^t\Sigma_{\bfX}C^{(2)} \;=\; \left(C^{(2)}\right)^t\Sigma_{\bfX}C^{(1)},
\]
and therefore the second equality of (7.24) is proven.

By letting $C^{(1)} = C^{(2)} = C = (c_i)_{i=1:n}$ in (7.24), we obtain that
\[
\var\!\left(\sum_{i=1}^n c_iX_i\right) \;=\; \cov\!\left(\sum_{i=1}^n c_iX_i,\sum_{i=1}^n c_iX_i\right) \;=\; C^t\Sigma_{\bfX}C.
\]
\end{proof}

The formulas from Lemma 7.3 connecting the covariance and correlation matrix of
random variables, as well as the Linear Transformation Property from section 7.3,
can also be established by using matrix operations and notations for multivariate
random variables. For more details, see section 7.7.

\begin{lemma}
(i) Any covariance matrix is symmetric positive semidefinite.

(ii) Any correlation matrix is symmetric positive semidefinite.
\end{lemma}

\begin{proof}
Let $\Sigma_{\bfX}$ and $\Omega_{\bfX}$ be the covariance and the correlation matrix
of the random variables $X_1, X_2, \ldots, X_n$, respectively. Recall from (7.8) and
(7.10) that $\Sigma_{\bfX}$ and $\Omega_{\bfX}$ are symmetric matrices.

(i) Let $v = (v_i)_{i=1:n} \in \R^n$ be an arbitrary column vector of size $n$. From
(7.25), we find that
\begin{equation}
v^t\Sigma_{\bfX}v \;=\; \var\!\left(\sum_{i=1}^n v_iX_i\right) \;\ge\; 0,
\end{equation}
since the variance of any random variable is nonnegative. Thus, $v^t\Sigma_{\bfX}v
\ge 0$ for any $v \in \R^n$, and we conclude that $\Sigma_{\bfX}$ is a symmetric
positive semidefinite matrix; cf. Definition 5.7.

(ii) Recall from (7.17) that
\begin{equation}
\Omega_{\bfX} \;=\; (D_{\sigma_{\bfX}})^{-1}\Sigma_{\bfX}(D_{\sigma_{\bfX}})^{-1},
\end{equation}
where $(D_{\sigma_{\bfX}})^{-1} = \diag\!\left(\frac{1}{\sigma_i}\right)_{i=1:n}$.

Let $v \in \R^n$ and let
\begin{equation}
w \;=\; (D_{\sigma_{\bfX}})^{-1}v.
\end{equation}
Then,
\begin{equation}
w^t \;=\; v^t\left((D_{\sigma_{\bfX}})^{-1}\right)^t \;=\; v^t(D_{\sigma_{\bfX}})^{-1},
\end{equation}
since $(D_{\sigma_{\bfX}})^{-1}$ is a diagonal matrix and therefore symmetric, i.e.,
$\left((D_{\sigma_{\bfX}})^{-1}\right)^t = (D_{\sigma_{\bfX}})^{-1}$.

From (7.29--7.32), we find that
\begin{align*}
v^t\Omega_{\bfX}v &= v^t(D_{\sigma_{\bfX}})^{-1}\Sigma_{\bfX}(D_{\sigma_{\bfX}})^{-1}v \\
&= w^t\Sigma_{\bfX}w \\
&\ge 0, \quad \forall\, w \in \R^n.
\end{align*}
Thus, $v^t\Omega_{\bfX}v \ge 0$ for all $v \in \R^n$, and we conclude that
$\Omega_{\bfX}$ is a symmetric positive semidefinite matrix; cf. Definition 5.7.
\end{proof}

\begin{lemma}
Let $X_1, X_2, \ldots, X_n$ be $n$ random variables on the same probability space.
The covariance matrix $\Sigma_{\bfX}$ of $X_1, X_2, \ldots, X_n$ is nonsingular if
and only if there is no linear combination of the random variables $X_1, X_2,
\ldots, X_n$ that would be constant, i.e., a constant random variable.
\end{lemma}

\begin{proof}
In order to prove Lemma 7.5, it is enough to show that the matrix $\Sigma_{\bfX}$
is singular if and only if there exist constants $c_1, c_2, \ldots, c_n$, and $c_0$,
not all equal to $0$, such that
\[
\sum_{i=1}^n c_iX_i + c_0 \;=\; 0,
\]
which is equivalent to
\begin{equation}
\var\!\left(\sum_{i=1}^n c_iX_i\right) \;=\; 0,
\end{equation}
since $c_0$ is a constant.

Let $C = (c_i)_{i=1:n}$. From (7.25) and (7.33), we obtain that
\begin{equation}
C^t\Sigma_{\bfX}C \;=\; 0.
\end{equation}
Note that, since the constants $c_1, c_2, \ldots, c_n$, and $c_0$ were assumed to
not be all equal to $0$, it follows that the $n \times 1$ vector $C$ is not equal to
$0$.

In other words, we need to show that
\begin{equation}
\Sigma_{\bfX} \text{ singular} \quad \Longleftrightarrow \quad \text{there exists } C \neq 0 \text{ such that } C^t\Sigma_{\bfX}C = 0.
\end{equation}
To prove (7.35), it is important to recall from Lemma 7.4 that $\Sigma_{\bfX}$ is a
symmetric positive semidefinite matrix.\footnote{The equivalence (7.35) does not
hold for a matrix which is not symmetric positive semidefinite matrix. For example,
if $A = \begin{pmatrix} 1 & 0 \\ 0 & -1 \end{pmatrix}$ and $C = \begin{pmatrix} 1 \\
1 \end{pmatrix}$, then $C^tAC = 0$ and $A$ is a nonsingular matrix.}

If $\Sigma_{\bfX}$ is a singular matrix, it follows from Theorem 4.3 that $0$ is an
eigenvalue of $\Sigma_{\bfX}$. Let $C \neq 0$ be an eigenvector of $\Sigma_{\bfX}$
corresponding to the eigenvalue $0$. Then, $\Sigma_{\bfX}C = 0$ and therefore
$C^t\Sigma_{\bfX}C = 0$.

If there exists a vector $C \neq 0$ such that $C^t\Sigma_{\bfX}C = 0$, we obtain
from Definition 5.6 that $\Sigma_{\bfX}$ is not symmetric positive definite. Then,
from Theorem 5.6, it follows that $0$ is an eigenvalue of $\Sigma_{\bfX}$, and we
conclude from Theorem 4.3 that the matrix $\Sigma_{\bfX}$ is singular.

Thus, the equivalence (7.35) is established.
\end{proof}
